\documentclass{article}
\usepackage{amsmath,amssymb}
\usepackage[margin=1in]{geometry}
\begin{document}
\title{Tail risk, execution, and allocation in the supplied double auction}
\author{Ada}
\date{}
\maketitle

\section*{Scope}
Q is \emph{Competitive Market Behavior of LLMs}. Its Sections 3.2--3.3 specify a
finite-horizon continuous double auction: crossing a standing quote executes
immediately at that quote, while an unfilled trader receives zero profit. Q's
Section 4.1 reports means across ten runs and links GPT Large's low trade count
to repeated small order improvements. The supplied P is \emph{Fixed-Batch CVaR
and Participation in a Quadratic Resource-Allocation Mechanism}. Its Sections
2--4 distinguish true from expected empirical CVaR and local from aggregate
risk. P's references to papers called Q and P in its introduction refer to
different papers from this assigned pair.

\section*{A conditional crossing result}
At a Q state, suppose a buyer can cross a standing ask below its reservation
price, giving certain profit \(r>0\). A seller has the symmetric option. Let
\(W\geq0\) be the eventual profit if the agent instead makes a noncrossing
improvement and follows some specified continuation policy. Write upper-tail
CVaR of loss \(-W\) with tail mass \(\beta\in(0,1]\). If
\[ \Pr(W=0)\geq\beta, \]
then \(\operatorname{CVaR}_{\beta}(-W)=0\), while immediate crossing has
\(\operatorname{CVaR}_{\beta}(-r)=-r\). Therefore minimizing local loss CVaR
strictly prefers crossing. This is a conditional decision result, not a claim
that Q's observed LLMs optimize CVaR. Q gives no counterfactual conditional
distribution of \(W\).
The buyer in Q does not observe other agents' reservation values, so a
fill-probability oracle would need to be specified or learned from allowed
public histories. The calculation does not give Q's traders this information.

For the transparent two-point case, let \(W=h>r\) with probability \(p\) and
\(W=0\) otherwise. Directly averaging the worst \(\beta\) probability mass
gives
\[
\operatorname{CVaR}_{\beta}(-W)=
\begin{cases}
0,&\beta\leq 1-p,\\
-h\bigl(\beta-(1-p)\bigr)/\beta,&\beta>1-p.
\end{cases}
\]
Crossing is preferred precisely when \(r>h\max\{0,1-(1-p)/\beta\}\), apart
from equality ties. Under the mean-profit criterion \(\beta=1\), the condition
is \(r>ph\). Thus \(ph>r>0\) and \(\beta\leq1-p\) give opposite choices.

\section*{Why local crossing does not prove market improvement}
Crossing a profitable standing quote creates positive immediate surplus, but
may use a seller who could later trade with a higher-value buyer. For example,
one seller has cost \(0.75\) and an ask of \(1\); a buyer valued at \(2.25\)
can cross now and obtain profit \(1.25\). If this buyer waits, a later buyer
valued at \(3.25\) may cross the same ask. On that path, the first trade creates
surplus \(1.50\) and the second creates \(2.50\). The values lie within Q's
reservation-price range. This is a pathwise displacement example, not a
counterfactual estimate for Q's full market. This is an allocation externality.
P's Section 4 gives a separate reason to distinguish local and market goals:
the sum of individual CVaRs generally differs from the CVaR of aggregate loss.

\section*{The missing aggregate tail}
Let \(E\in[0,1]\) be Q's normalized allocative efficiency in a run and
\(L=1-E\) its inefficiency loss. Q's round-five GPT Small mean \(E=0.91\)
does not identify \(\operatorname{CVaR}_{0.1}(L)\). The ten efficiency values
\((0.91,\ldots,0.91)\) and \((0.1,1,\ldots,1)\) have the same mean, but their
empirical upper-tenth loss CVaRs are \(0.09\) and \(0.9\), respectively. This
only shows insufficiency of the efficiency mean; Q reports other marginal
means, and the example does not claim to reproduce its joint data.

If a severe failure occurs independently with probability \(0.05\), ten
independent runs miss it with probability \(0.95^{10}\approx0.599\). At least
\[
\left\lceil\frac{\log(0.05)}{\log(0.95)}\right\rceil=59
\]
independent runs are needed
for a probability of at least \(0.95\) of seeing one such failure. This is an
event-detection calculation, not a CVaR confidence interval. P's Section 4
bound \(e_n=L_n\delta_n+U_n\sqrt{2\pi}/(\beta_n\sqrt{s_n})\) concerns an
expected empirical surrogate and likewise does not certify a realized batch.

\section*{Research test and limits}
A focused experiment could randomize an explicitly local CVaR-informed
crossing rule against Q's existing trader prompts at matched market states.
It would measure conditional nonfill rates, valuation ranks of trading agents,
mean surplus, and the tail of market inefficiency across independent runs. The
testable question is whether fewer missed profitable trades outweigh displaced
high-value traders. Sampling and reporting must address the aggregate tail
separately from individual risk. This is a proposed experiment, not a finding
about the intervention. Existing risk-based CDA bidding predates both supplied
papers: Vytelingum et al. (2004),
\texttt{https://eprints.soton.ac.uk/259567/}. CVaR also appears in
double-auctioned power-market bidding (Panda et al., 2019,
\texttt{https://doi.org/10.1049/iet-gtd.2018.5905}). The decision threshold
above is elementary and is not a standalone novelty claim. A possible
contribution would require the controlled LLM-market test and aggregate-risk
measurement described here.
P's equilibrium and optimization guarantees do not directly transfer to Q:
the crossing decision is discrete, Q's matching rule is discontinuous at the
spread, and Q has neither P's quadratic taxes nor its empirical-CVaR oracle.
Q's ten evaluation runs are separate from any oracle batch size used by a
trader.
\end{document}
